% ============================================================
% 5.1 Hiện thực: Mua hàng
% ============================================================

\section{Hiện thực --- Mua hàng}
\label{sec:impl-order}

\subsection{Mô tả cách sử dụng}

Quy trình mua hàng trên ShopNexus gồm hai giai đoạn chính. Trong giai đoạn \textbf{checkout}, người mua chọn sản phẩm, địa chỉ giao hàng và đơn vị vận chuyển cho từng item, sau đó thanh toán ngay (kết hợp ví nội bộ và cổng thanh toán). Hệ thống đặt chỗ tồn kho, snapshot giá, và tạo các mục chờ liên kết với phiên thanh toán. Trong giai đoạn \textbf{xác nhận}, người bán xem danh sách các mục đã thanh toán, gộp và xác nhận hoặc từ chối. Khi xác nhận, hệ thống tạo vận đơn và đơn hàng chính thức. Nếu người bán không phản hồi trong 48 giờ, hệ thống tự động hủy và hoàn tiền vào ví.

\begin{figure}[!htbp]
\centering
\includegraphics[width=\textwidth,keepaspectratio]{uploads/viewcart.png}
\caption{Trang giỏ hàng của người mua}
\label{fig:view-cart-page}
\end{figure}

Trang giỏ hàng hiển thị danh sách sản phẩm đã chọn, cho phép người mua điều chỉnh số lượng hoặc xóa sản phẩm trước khi chuyển sang trang checkout.

\begin{figure}[!htbp]
\centering
\includegraphics[width=\textwidth,keepaspectratio]{uploads/checkout.png}
\captionof{figure}{Trang checkout và thanh toán}
\label{fig:impl-checkout}
\end{figure}

Trang checkout cho phép buyer chọn địa chỉ giao hàng, đơn vị vận chuyển cho từng item, và phương thức thanh toán. Sau khi thanh toán, hệ thống tạo các mục chờ seller xác nhận.

\begin{figure}[!htbp]
\centering
\includegraphics[width=\textwidth,keepaspectratio]{uploads/buyervieworders.png}
\captionof{figure}{Danh sách đơn hàng của người mua}
\label{fig:impl-order-list}
\end{figure}


Trang đơn hàng cho phép buyer theo dõi trạng thái và thực hiện các thao tác như hủy đơn hoặc yêu cầu hoàn trả.

% ------------------------------------------------------------------
\subsection{Đánh giá}
\label{subsec:order-eval}

\small
\rowcolors{2}{white}{rowgray}
\begin{longtable}{@{}L{0.5cm}L{6.5cm}L{0.5cm}L{6.5cm}@{}}
\caption{Đánh giá chức năng mua hàng}\label{tab:order-eval}\\
\toprule
\rowcolor{headerblue}
\multicolumn{2}{c}{\thead{Điểm mạnh}} & \multicolumn{2}{c}{\thead{Hạn chế}} \\
\midrule
\endfirsthead
\toprule
\rowcolor{headerblue}
\multicolumn{2}{c}{\thead{Điểm mạnh}} & \multicolumn{2}{c}{\thead{Hạn chế}} \\
\midrule
\endhead
\midrule
\multicolumn{2}{r}{\small\textit{Tiếp theo trang sau}} \\
\endfoot
\bottomrule
\endlastfoot

1 & Reserve inventory tại thời điểm checkout, tránh oversell & 1 & Chưa hỗ trợ thanh toán quốc tế (Stripe, PayPal) \\
2 & Checkout hai bước (pending $\rightarrow$ confirm) cho phép seller kiểm soát & 2 & Chưa có push notification cho mobile \\
3 & Tích hợp VNPay với nhiều phương thức (QR, Bank, ATM, COD) & 3 & Chưa hỗ trợ mua từ nhiều seller trong một checkout \\
4 & Durable execution (Restate) đảm bảo exactly-once cho luồng thanh toán & 4 & Chưa có theo dõi vận chuyển real-time \\
\end{longtable}
\rowcolors{1}{}{}
\normalsize

% ------------------------------------------------------------------
\subsection{Kiểm thử hộp đen}
\label{subsec:order-blackbox}

Kiểm thử hộp đen cho chức năng mua hàng tập trung vào hành vi quan sát được từ phía người dùng, không quan tâm đến cấu trúc mã bên trong. Test case được thiết kế theo kỹ thuật \textbf{phân lớp tương đương} (equivalence partitioning) --- mỗi kịch bản đại diện cho một lớp đầu vào --- kết hợp \textbf{phân tích giá trị biên} (boundary value analysis) cho các trường hợp giới hạn (số lượng = 0, tồn kho vừa đủ, hết hạn thanh toán).

\small
\rowcolors{2}{white}{rowgray}
\begin{longtable}{@{}L{1.2cm}L{3cm}L{4cm}L{4.5cm}L{0.8cm}@{}}
\caption{Test case hộp đen --- chức năng mua hàng}\label{tab:order-test}\\
\toprule
\rowcolor{headerblue}
\thead{ID} & \thead{Kịch bản} & \thead{Điều kiện trước / Đầu vào} & \thead{Kết quả mong đợi} & \thead{KQ} \\
\midrule
\endfirsthead
\toprule
\rowcolor{headerblue}
\thead{ID} & \thead{Kịch bản} & \thead{Điều kiện trước / Đầu vào} & \thead{Kết quả mong đợi} & \thead{KQ} \\
\midrule
\endhead
\midrule
\multicolumn{5}{r}{\small\textit{Tiếp theo trang sau}} \\
\endfoot
\bottomrule
\endlastfoot
TC-01 & Thêm sản phẩm vào giỏ hàng & Giỏ trống; SKU hợp lệ, qty = 1 & Giỏ hàng có 1 item đúng SKU và số lượng & Đạt \\
TC-02 & Thêm sản phẩm hết hàng & SKU tồn kho = 0 & Hiển thị lỗi ``Sản phẩm hết hàng'' & Đạt \\
TC-03 & Thêm trùng SKU đã có trong giỏ & Giỏ có SKU A qty = 1; thêm lại SKU A qty = 2 & Qty cập nhật thành 3, không tạo dòng mới & Đạt \\
TC-04 & Checkout tạo pending items & Giỏ hàng có 2 items, chọn địa chỉ và vận chuyển & Items Pending được tạo, inventory reserved & Đạt \\
TC-05 & Checkout với qty vượt tồn kho & SKU tồn kho = 5, qty = 10 & Hệ thống từ chối, hiển thị lỗi tồn kho & Đạt \\
TC-06 & Buyer hủy pending item & Item ở trạng thái chờ xác nhận & Item hủy, inventory hoàn trả, ví hoàn tiền & Đạt \\
TC-07 & Seller xác nhận pending items & Items cùng buyer, cùng địa chỉ, cùng vận chuyển & Order + transport tạo, items gắn order\_id & Đạt \\
TC-08 & Seller từ chối pending items & Items ở trạng thái chờ xác nhận & Items hủy, ví buyer hoàn tiền đầy đủ & Đạt \\
TC-09 & Thanh toán thành công (webhook) & Phiên thanh toán Pending, webhook trả Success & Payment $\to$ Success, lên lịch timeout 48h & Đạt \\
TC-10 & Thanh toán thất bại (webhook) & Phiên thanh toán Pending, webhook trả Failed & Payment $\to$ Failed, items hủy, kho hoàn trả & Đạt \\
TC-11 & Checkout ví + cổng thanh toán & Số dư ví $<$ tổng tiền & Ví trừ một phần, phiên cổng tạo cho phần còn lại & Đạt \\
TC-12 & Xem đơn hàng với bộ lọc & Buyer có 5 đơn, lọc theo trạng thái & Danh sách trả về đúng theo bộ lọc & Đạt \\
TC-13 & Mua ngay (BuyNow) & SKU hợp lệ, qty = 1 & Tạo item không xóa giỏ hàng hiện tại & Đạt \\
TC-14 & Mua ngay nhiều SKU & BuyNow = true, 2 SKUs & Bị từ chối: BuyNow chỉ cho phép 1 SKU & Đạt \\
\end{longtable}
\rowcolors{1}{}{}
\normalsize
